\chapter{绪论}

电力电子技术是电气工程学科中最活跃的分支之一。它研究的是“如何用半导体开关把电能从一种形态搬运成另一种形态”，因此它既要有电路理论作骨架，又要有控制理论作神经。本章回答四个问题：什么是电力电子技术、它是怎么发展起来的、它用在哪里、这本教材怎么读。

% ============================================================
\section{什么是电力电子技术}

\subsection{定义与学科定位}

\begin{tcolorbox}[enhanced,breakable,colback=highlight,colframe=yhbAccent!30!white,boxrule=0.5pt,arc=2.5pt,left=11pt,right=11pt,top=9pt,bottom=9pt,boxsep=0pt,before skip=11pt,after skip=11pt]
\textbf{\color{yhbAccent}定义（电力电子技术）}\quad 电力电子技术是\textbf{使用电力电子器件对电能进行变换和控制的技术}。目前所用的电力电子器件均由半导体制成，故也称电力半导体器件。
\end{tcolorbox}

几点必须分清：

\begin{itemize}
  \item 电子技术分为\textbf{信息电子技术}与\textbf{电力电子技术}两大分支。通常所说的模拟电子技术、数字电子技术都属于信息电子技术。
  \item 信息电子技术主要用于\textbf{信息处理}，电力电子技术主要用于\textbf{电力变换}——这是二者本质上的不同。
  \item 电力电子技术中所变换的“电力”功率可以大到数百兆瓦甚至吉瓦，也可以小到数瓦甚至毫瓦级，跨度极大。
  \item 广义而言，电子电路中的功率放大和功率输出部分也可算作电力电子电路；各种电子装置（电视机、计算机）的电源部分都是电力电子电路。
\end{itemize}

\subsection{电力变换的四大类}

常用的电力有交流和直流两种。从公用电网直接得到的是交流，从蓄电池和干电池得到的是直流。这些电源往往不能直接满足要求，需要进行电力变换。电力变换通常分为四大类，见表~\ref{tab:1-1}。

\begin{table}[htbp]
\centering
\caption{电力变换的种类}
\label{tab:1-1}
\begin{tabular}{llll}
\toprule
\textbf{变换类型} & \textbf{英文缩写} & \textbf{习惯名称} & \textbf{实现电路} \\
\midrule
交流 $\rightarrow$ 直流 & AC--DC & 整流（Rectifier） & 整流电路 \\
直流 $\rightarrow$ 交流 & DC--AC & 逆变（Inverter） & 逆变电路 \\
直流 $\rightarrow$ 直流 & DC--DC & 直流斩波 & 直流斩波电路 \\
交流 $\rightarrow$ 交流 & AC--AC & 交流电力控制 / 变频 / 变相 & 交交变频、交流调压 \\
\bottomrule
\end{tabular}
\end{table}

\corepoint{直流变直流\textbf{并非电力种类的变换}，而是电压（或电流）的变换，即一种直流电压（或电流）变为另一种直流电压（或电流）；交流变交流除了电压或电流的变换外，还多了\textbf{频率或相数}的变换。进行上述电力变换的技术统称\textbf{变流技术}。}

\subsection{两个分支：器件制造技术与变流技术}

通常把电力电子技术分为\textbf{电力电子器件制造技术}和\textbf{变流技术}两个分支。变流技术也称为电力电子器件的应用技术，包括用电力电子器件构成各种电力变换电路、对这些电路进行控制、以及由这些电路构成电力电子装置和系统的技术。

\begin{itemize}
  \item 器件制造技术的理论基础是\textbf{半导体物理}；
  \item 变流技术的理论基础是\textbf{电路理论}；
  \item 没有晶闸管、电力晶体管、IGBT 等器件，就没有电力电子技术，所以\textbf{器件制造技术是基础}；而电力电子技术主要用于电力变换，所以\textbf{变流技术是核心}。
\end{itemize}

\subsection{“倒三角形”：电力电子学的学科交叉}

电力电子学（Power Electronics）这一名称出现在 20 世纪 60 年代（比晶闸管的出现晚）。1974 年，美国学者 W. Newell 用倒三角形描述了这一学科：它由电力学、电子学和控制理论三个学科交叉形成，这一观点被全世界普遍接受。

\begin{figure}[htbp]
\centering
\begin{tikzpicture}[font=\small]
  \draw[yhbAccent,thick] (0,0) -- (6,0) -- (3,-3.2) -- cycle;
  \node[text=yhbAccent] at (3,-1.55) {\textbf{电力电子技术}};
  \node at (1.55,-0.45) {电力学};
  \node at (4.45,-0.45) {电子学};
  \node at (3,-2.85) {控制理论};
  \node[text=textgray,font=\footnotesize] at (3,0.35) {连续、离散控制理论};
\end{tikzpicture}
\caption{描述电力电子学的倒三角形}
\label{fig:1-1}
\end{figure}

“电力电子学”和“电力电子技术”是分别从学术和工程技术两个不同角度来称呼的，其实际内容并没有很大的不同。

\subsection{与信息电子技术的根本区别}

电力电子技术与电子学的关系是显见的：信息电子学可分为电子器件和电子电路两大分支，这分别与电力电子器件和电力电子电路相对应；二者的制造理论基础（半导体理论）相同，大多数工艺也相同，现代电力电子器件的制造大都使用集成电路制造工艺。

差别在应用目的与工作状态：

\corepoint{信息电子技术中，半导体器件既可处于放大状态，也可处于开关状态；而电力电子技术中，\textbf{为避免功率损耗过大，电力电子器件总是工作在开关状态}。这是电力电子技术区别于信息电子技术的一个重要特征。}

电力电子技术广泛用于电气工程中，这是电力电子学和电力学的主要关系。电力电子装置广泛应用于高压直流输电、静止无功补偿、电力机车牵引、交直流电力传动、电解、励磁、电加热、高性能交直流电源之中，因此国内外通常把电力电子技术归属于电气工程学科。在我国的学科分类中，电气工程是一级学科，包含五个二级学科：电力系统及其自动化、电机与电器、高电压与绝缘技术、电力电子与电力传动、电工理论与新技术。

% ============================================================
\section{电力电子技术的发展史}

电力电子器件的发展对电力电子技术的发展起着决定性的作用，因此电力电子技术的发展史是以电力电子器件的发展史为纲的。

\subsection{史前期（黎明期）}

\begin{itemize}
  \item 1904 年出现\textbf{电子管}，它能在真空中对电子流进行控制，并应用于通信和无线电，从而开启了电子技术用于电力领域的先河。
  \item 后来出现\textbf{水银整流器}（汞弧整流器），把水银封于管内，利用对其蒸气的点弧可对大电流进行控制，其性能和晶闸管已经非常相似。缺点是水银对人体有害、电压降落也很高。
  \item 20 世纪 30--50 年代是水银整流器发展迅速并大量应用的时期，广泛用于电化学工业、电气铁道直流变电所、轧钢用直流电动机传动，甚至用于直流输电。这一时期，各种整流电路、逆变电路、周波变流电路的理论已经发展成熟并广为应用。
  \item 把交流变为直流的方法，除水银整流器外，还有发展更早的\textbf{电动机—直流发电机组}，即\textbf{变流机组}（旋转变流机组）。与旋转变流机组相对应，静止变流器的称呼从水银整流器开始而沿用。
\end{itemize}

\subsection{晶闸管时代}

1947 年，美国贝尔实验室发明了\textbf{晶体管}，引发了电子技术的一场革命；最先用于电力领域的半导体器件是硅二极管。1957 年，美国通用电气公司研制出第一个\textbf{晶闸管}，一般以此作为电力电子技术诞生的标志。

\corepoint{晶闸管是\textbf{通过对门极的控制能够使其导通而不能使其关断}的器件，属于\textbf{半控型器件}；对晶闸管电路的控制方式主要是\textbf{相位控制方式}（简称相控方式）；晶闸管的关断通常依靠电网电压等外部条件来实现。这使晶闸管的应用受到很大局限。}

晶闸管出现后，由于其优越的电气性能和控制性能，很快取代了水银整流器和旋转变流机组，应用范围迅速扩大。电化学工业、铁道电气机车、钢铁工业（轧钢电气传动、感应加热等）、电力工业（直流输电、无功补偿等）的迅速发展也给晶闸管提供了用武之地。\textbf{电力电子技术的概念和基础，就是由于晶闸管及晶闸管变流技术的发展而确立的}。

\subsection{全控型器件与复合型器件}

20 世纪 70 年代后期，以\textbf{门极可关断晶闸管（GTO）}、\textbf{电力双极型晶体管（BJT）}和\textbf{电力场效应晶体管（Power MOSFET）}为代表的全控型器件迅速发展。

\begin{itemize}
  \item 全控型器件的特点是：通过对门极（基极、栅极）的控制\textbf{既可使它开通又可使它关断}；开关速度普遍高于晶闸管，可用于开关频率较高的电路。
  \item 与晶闸管电路的相控方式相对应，采用全控型器件的电路主要控制方式为\textbf{脉冲宽度调制（PWM）}方式，相对于相位控制方式可称为\textbf{斩波控制方式}（斩控方式）。
\end{itemize}

20 世纪 80 年代后期，以\textbf{绝缘栅极双极型晶体管（IGBT）}为代表的\textbf{复合型器件}异军突起。IGBT 是 MOSFET 和 BJT 的复合，它把 MOSFET 驱动功率小、开关速度快的优点，和 BJT 通态压降小、载流能力大、可承受电压高的优点集于一身，性能十分优越，成为现代电力电子技术的主导器件。与之对应，\textbf{MOS 控制晶闸管（MCT）}和\textbf{集成门极换流晶闸管（IGCT）}都是 MOSFET 和 GTO 的复合，也综合了两种器件的优点，其中 IGCT 已经获得大量应用。

为使装置结构紧凑、体积减小，常把若干个电力电子器件及必要的辅助元件做成\textbf{模块}形式；后来又进一步把驱动、控制、保护电路和器件集成在一起，构成\textbf{电力电子集成电路（PIC）}。目前 PIC 的功率还较小、电压也较低，它面临着电压隔离（主电路为高压而控制电路为低压）、热隔离（主电路发热严重）、电磁干扰（开关器件通断高压大电流且与控制电路同处一个芯片）等几大难题，但这代表了电力电子技术发展的一个重要方向。电力电子集成技术的焦点是\textbf{混合集成技术}，即把不同的单个芯片集成封装在一起；此外\textbf{系统集成}也是重要方面，对超大功率集成技术尤其如此。

随着全控型器件不断进步，电力电子电路的工作频率不断提高，开关损耗也随之增大。为了减小开关损耗，\textbf{软开关技术}应运而生——\textbf{零电压开关（ZVS）}和\textbf{零电流开关（ZCS）}就是软开关的最基本形式。理论上讲采用软开关技术可使开关损耗降为零，可以提高效率，也使开关频率得以进一步提高，从而提高装置的功率密度。

\begin{table}[htbp]
\centering
\caption{电力电子技术发展阶段的粗略对照}
\label{tab:1-2}
\begin{tabular}{p{2.6cm}p{4.2cm}p{6.4cm}}
\toprule
\textbf{阶段} & \textbf{代表器件} & \textbf{技术特征} \\
\midrule
史前期（$\sim$1957 前） & 电子管、水银整流器、变流机组 & 旋转变流、静止汞弧整流；整流/逆变/周波变流理论已成熟 \\
晶闸管时代（1957 起） & 晶闸管（半控型） & 相控方式；关断依赖电网等外部条件；装置容量大、频率低 \\
全控型时代（70 年代后期起） & GTO、BJT、Power MOSFET & 门极可开通亦可关断；斩控（PWM）方式；开关频率提高 \\
复合型时代（80 年代后期起） & IGBT、MCT、IGCT & MOS 栅极 + 双极导电，驱动功率小且通态压降小；成为主导器件 \\
集成化与软开关 & 功率模块、PIC、ZVS/ZCS & 结构紧凑、功率密度高；开关损耗趋近于零 \\
\bottomrule
\end{tabular}
\end{table}

% ============================================================
\section{电力电子技术的应用}

电力电子技术的应用范围十分广泛。它不仅用于一般工业，也广泛用于交通运输、电力系统、通信系统、计算机系统、新能源系统等，在照明、空调等家用电器中也有广泛应用。

\subsection{一般工业}

工业中大量应用各种交直流电动机。直流电动机有良好的调速性能，为其供电的可控整流电源或直流斩波电源都是电力电子装置。近年来由于电力电子变频技术的迅速发展，交流电动机的调速性能可与直流电动机相媲美，交流调速技术逐渐大量应用并占据了主导地位。大至数千千瓦的轧钢机，小到几百瓦的数控机床伺服电动机，以及矿山牵引等场合，都广泛采用电力电子交流调速技术。一些调速性能要求不高的大型鼓风机也采用了变频装置以达到节能目的。还有些并不特别要求调速的电动机，为避免起动时的电流冲击而采用\textbf{软起动装置}，这也是电力电子装置。

\corepoint{电动机所消耗的电力甚至达到发电厂所发电力的 60\% 以上，以至于有人认为\textbf{电力传动是电力电子技术的“主战场”}。}

电化学工业大量使用直流电源，电解铝、电解食盐水等都需要大容量整流电源，电镀装置也需要整流电源。电力电子技术还大量用于冶金工业中的高频或中频感应加热电源、淬火电源及直流电弧炉电源等场合。

\subsection{交通运输}

电气化铁道中广泛采用电力电子技术：电气机车中的直流机车采用整流装置，交流机车采用变频装置；直流斩波器也广泛用于铁道车辆。在磁悬浮列车中，电力电子技术更是一项关键技术。除牵引电动机传动外，车辆中的各种辅助电源也都离不开电力电子技术。电动汽车的电动机依靠电力电子装置进行电力变换和驱动控制，蓄电池的充电也离不开电力电子装置。一台高级汽车中需要许多控制电动机，它们也要靠变频器和斩波器驱动并控制。航空和航海都离不开电力电子技术（电源和驱动）；电梯若算作交通运输，以前的电梯大都采用直流调速系统，近年来交流变频调速已成为主流。

\subsection{电力系统}

电力电子技术在电力系统中有着非常广泛的应用。据估计，发达国家在用户最终使用的电能中，有 60\% 以上的电能至少经过一次以上电力电子变流装置的处理。可以毫不夸张地说，如果离开电力电子技术，电力系统的现代化就是不可想象的。

\begin{itemize}
  \item \textbf{直流输电}：长距离、大容量输电时有很大优势，其送电端的整流阀和受电端的逆变阀都采用晶闸管变流装置；轻型直流输电则主要采用全控型的 IGBT 器件。
  \item \textbf{柔性交流输电（FACTS）}：依靠电力电子装置才得以实现。
  \item \textbf{无功补偿与谐波抑制}：晶闸管控制电抗器（TCR）、晶闸管投切电容器（TSC）都是重要的无功补偿装置；采用全控型器件（如 IGBT）的静止无功发生器（SVG）、有源电力滤波器（APF）等新型装置性能更为优越。
  \item \textbf{电能质量控制}：在配电网中，电力电子装置还可用于防止电网瞬时停电、瞬时电压跌落、闪变等，以改善供电质量。
  \item \textbf{变电所}：给操作系统提供可靠的交直流操作电源、给蓄电池充电等，都需要电力电子装置。
\end{itemize}

\subsection{电子装置用电源}

各种电子装置一般都需要不同电压等级的直流电源供电。通信设备中的程控交换机所用直流电源以前用晶闸管整流电源，现在已改为采用全控型器件的高频开关电源；大型计算机所需的工作电源、微型计算机内部的电源现在也都采用高频开关电源。以前大量采用的线性稳压电源，由于高频开关电源体积小、重量轻、效率高，现已逐渐被取代。可以说，信息技术装置都需要电力电子装置提供电源，\textbf{信息电子技术离不开电力电子技术}。在有大型计算机等场合，常常需要\textbf{不间断电源（UPS）}供电。

\subsection{家用电器与其他}

照明在家用电器中占有十分突出的地位。电力电子照明电源体积小、发光效率高、可节省大量能源，通常采用电力电子装置的光源被称为“节能灯”，正在逐步取代传统的白炽灯和荧光灯。\textbf{变频空调器}是家用电器中应用电力电子技术的典型例子；电视机、音响设备、家用计算机等电子设备的电源部分也都需要电力电子技术；不少洗衣机、电冰箱、微波炉中也应用了电力电子技术。

新能源及新型发电方式越来越受到重视，其中风力发电、太阳能发电发展最为迅速，燃料电池备受关注。太阳能发电和风力发电受环境的制约，发出的电力质量较差，常需要储能装置缓冲、需要改善电能质量，这就需要电力电子技术；当需要和电力系统联网时，更离不开电力电子技术。抽水储能发电站中大型电动机的起动和调速、超导储能的强大直流电源、核聚变反应堆的大容量脉冲电源，以及科学实验中的各种特种电源，都是电力电子技术的用武之地。

\corepoint{电力电子装置提供给负载的是各种不同的直流电源、恒频交流电源以及变频交流电源，因此也可以说，\textbf{电力电子技术研究的就是电源技术}。它在大型风机、水泵变频调速和照明电源方面的节能效果十分显著，因此也被称为\textbf{节能技术}。以前电力电子技术的应用偏重于中、大功率，现在在 1 kW 以下甚至几十瓦以下的功率范围内应用越来越广，这已成为一个重要的发展趋势。}

% ============================================================
\section{本教材的内容简介和使用说明}

本教材除第 1 章绪论外，可分为四大部分。

\begin{table}[htbp]
\centering
\caption{本教材的内容结构}
\label{tab:1-3}
\begin{tabular}{p{3.4cm}p{2.2cm}p{8.0cm}}
\toprule
\textbf{部分} & \textbf{章} & \textbf{内容与地位} \\
\midrule
器件 & 第 2 章、第 9 章 & 第 2 章是全书的基础，介绍各种器件的基本结构、工作原理、主要参数、应用特性，以晶闸管、IGBT、电力 MOSFET 三种应用最广的器件为重点；第 9 章是器件应用的共性问题，包括驱动、控制、保护以及串并联 \\
电路 & 第 3--6 章 & 各种基本电力变换电路：整流、逆变、直流—直流变流、交流—交流变流 \\
控制技术 & 第 7 章、第 8 章 & PWM 控制技术（没有新的电路拓扑，但使电路性能大为改善，适用于表~\ref{tab:1-1} 所列各类变换电路）；软开关技术（一般并不改变原来的基本电路，而是附加一些电路以实现软开关） \\
应用 & 第 10 章 & 电力电子装置的应用，体现第 5 版与第 4 版的重要差异 \\
\bottomrule
\end{tabular}
\end{table}

与前三版教材相比，第 4 版完成了以晶闸管电路为主体内容向以采用全控型器件的电路为主体内容的转移；第 5 版则在第 4 版问世 8 年来的新进展基础上做了充分反映。

学习建议：

\begin{itemize}
  \item 本课程要求先修“电路”和“电子技术基础”，最好也学过“自动控制理论”，并能较熟练地掌握示波器等电子仪器的使用方法。
  \item 电力电子技术实践性很强，实验在教学中有重要位置；教材正文后附“教学实验”部分，精选了五个最基本、有较高实用价值的实验。
  \item 术语在一门课程的学习中起着“纲”的作用，书末附有中英文“术语索引”，遇到拿不准的概念应回头查定义。
  \item 各章前有前言说明本章的地位及概要，大部分章末有小结，对全章要点和重点进行总结——复习时应以“小结核对 + 公式自推”为主，而不是背结论。
\end{itemize}

\corepoint{本章的考点集中在\textbf{概念辨析}：四大类电力变换的名称与英文缩写、半控型与全控型器件的分界、相控方式与斩控（PWM）方式的对应关系、“器件制造技术是基础、变流技术是核心”这一论断，以及电力电子器件\textbf{总工作在开关状态}这一与信息电子技术的根本区别。}
